\begingroup
% Distinguish source identifiers and permit wrapping without changing their text.
\def\UrlFont{\small\ttfamily}
\setlength{\LTpre}{0.5\baselineskip}
\setlength{\LTpost}{0.5\baselineskip}
\section{End-to-end Lean verification and reproduction guide}
\label{app:lean-verification}

This appendix describes the accompanying Lean development, identifies the
precise statements that it proves, and gives instructions for rebuilding and
inspecting the proofs.  The author of the development is Karthik Sheshadri.  The sources and the
archived release, which also contains this paper, are at
\begin{itemize}[leftmargin=*]
\item Lean sources: \leanRepositoryURL;
\item archived release: \leanArchiveDOI.
\end{itemize}
The entry point is \nolinkurl{OriginalConormal.lean}; the final numerical
theorems are in \nolinkurl{Sheshadri/ConormalExactMultiplicity.lean} and
\nolinkurl{Sheshadri/ConormalBorderMultiplicity.lean}.

\begin{readerbox}{What ``end to end'' means here}
For every pair of natural numbers $n,m$ with $n\geq3$, the development proves
the ordinary and symmetric lower bounds from a literal determinant
representation, or from literal border membership, respectively.  The final
theorems do not ask the user to provide a B\'ezout theorem, a generic flag,
a degree formula, a finite section, a specialization principle, or a geometric
certificate.  These inputs are proved or constructed in the dependency chain.

The formal implementation uses \emph{local Gauss sections and their full
intersection lengths} to obtain the numerical consequences of
conormal/Gauss specialization.  It is not a claim that every intermediate
global-cycle assertion in the main text has been formalized verbatim.
Section~\ref{app:lean-scope} specifies this distinction.
\end{readerbox}

\subsection{The statements checked by Lean}
\label{app:lean-statements}

Throughout the development, the target is the actual multivariate polynomial
\[
 f_n=\sum_{i\in\{0,\ldots,n-1\}}X_i^n
 \quad\text{in}\quad \CC[X_0,\ldots,X_{n-1}].
\]
This indexing differs from $F_n=x_1^n+\cdots+x_n^n$ only by renaming variables.
An affine matrix pencil is an $(n+1)$-tuple of constant $m\times m$ matrices,
and its polynomial is literally
\[
 \det\left(A_0+\sum_{i=0}^{n-1}X_iA_{i+1}\right).
\]
In the symmetric model every coefficient matrix is required to equal its
transpose.  The definitions occur in \nolinkurl{PolynomialModel.lean}, in
the namespace \nolinkurl{Sheshadri.PolynomialModel}.

Put
\begin{align*}
 C(n)&=n(n-1)^{n-2},\\
 B(m,n)&=\sum_{i=1}^{n-1}
   \binom mi\binom{m-1}{n-1-i}\binom{n-2}{i-1},\\
 B_{\mathrm{sym}}(m,n)&=2^{n-2}\binom m{n-1}.
\end{align*}
Their Lean names are, respectively, \nolinkurl{fermatCount},
\nolinkurl{ordinaryBezout}, and \nolinkurl{symmetricBezout}.  These are
natural-number definitions, including the usual zero values of binomial
coefficients outside their range.  The final real-number inequalities use
\texttt{\detokenize{Real.exp 1}} for $e$.

The following is a complete account of the inputs of the final endpoints.
All entries quantify only $n,m\in\mathbb N$, have the additional hypothesis
$3\leq n$, and have no further implicit geometric or typeclass parameter.
The short names in the first table are in the namespace
\nolinkurl{Sheshadri.ConormalExactMultiplicity}.

\begingroup\small
\begin{longtable}{@{}>{\raggedright\arraybackslash}p{0.36\textwidth}>{\raggedright\arraybackslash}p{0.40\textwidth}>{\raggedright\arraybackslash}p{0.18\textwidth}@{}}
\toprule
Lean theorem & Remaining representation hypothesis & Conclusion\\
\midrule
\endhead
\nolinkurl{ordinary_exact_count} &
\nolinkurl{ExactRepresentation}\newline\texttt{n m Complex} & $C(n)\leq B(m,n)$\\[2pt]
\nolinkurl{symmetric_exact_count} &
\nolinkurl{SymmetricExactRepresentation}\newline\texttt{n m Complex} &
$C(n)\leq B_{\mathrm{sym}}(m,n)$\\[2pt]
\nolinkurl{ordinary_exact_quadratic} &
\nolinkurl{ExactRepresentation}\newline\texttt{n m Complex} &
$\displaystyle\frac{(n-1)^2}{4e}\leq m$\\[5pt]
\nolinkurl{symmetric_exact_quadratic} &
\nolinkurl{SymmetricExactRepresentation}\newline\texttt{n m Complex} &
$\displaystyle\frac{(n-1)^2}{2e}\leq m$\\
\bottomrule
\end{longtable}
\endgroup

The border names below are in
\nolinkurl{Sheshadri.ConormalBorderMultiplicity}.  In each paired row,
``ordinary'' and ``symmetric'' select the corresponding hypothesis and
conclusion.

\begingroup\small
\begin{longtable}{@{}>{\raggedright\arraybackslash}p{0.38\textwidth}>{\raggedright\arraybackslash}p{0.38\textwidth}>{\raggedright\arraybackslash}p{0.18\textwidth}@{}}
\toprule
Lean theorem & Remaining border hypothesis & Conclusion\\
\midrule
\endhead
\nolinkurl{ordinary_border_count},
\nolinkurl{symmetric_border_count} &
\nolinkurl{OrdinaryFiniteSupportBorder} \texttt{n m},
\nolinkurl{SymmetricFiniteSupportBorder} \texttt{n m} &
$C(n)\leq B(m,n)$ or $B_{\mathrm{sym}}(m,n)$\\[4pt]
\nolinkurl{ordinary_border_quadratic},
\nolinkurl{symmetric_border_quadratic} &
The same finite-support predicates &
$\displaystyle\frac{(n-1)^2}{4e}\leq m$ or
$\displaystyle\frac{(n-1)^2}{2e}\leq m$\\[7pt]
\nolinkurl{ordinary_borderAtSize_count},
\nolinkurl{symmetric_borderAtSize_count} &
\nolinkurl{OrdinaryBorderAtSize} \texttt{n m},
\nolinkurl{SymmetricBorderAtSize} \texttt{n m} &
The corresponding count bound\\[4pt]
\nolinkurl{ordinary_borderAtSize_quadratic},
\nolinkurl{symmetric_borderAtSize_quadratic} &
The same canonical-support predicates &
The corresponding quadratic bound\\[4pt]
\nolinkurl{ordinary_zariski_count},
\nolinkurl{symmetric_zariski_count} &
\nolinkurl{OrdinaryZariskiBorder} \texttt{n m},
\nolinkurl{SymmetricZariskiBorder} \texttt{n m} &
The corresponding count bound\\[4pt]
\nolinkurl{ordinary_zariski_quadratic},
\nolinkurl{symmetric_zariski_quadratic} &
The same Zariski-border predicates &
The corresponding quadratic bound\\
\bottomrule
\end{longtable}
\endgroup

For example, the full mathematical content of the ordinary finite-support
border endpoint is exactly
\[
 \forall n,m\in\mathbb N,\quad
 3\leq n\ \Longrightarrow\
 \operatorname{OrdinaryFiniteSupportBorder}(n,m)
 \ \Longrightarrow\
 \frac{((n:\mathbb R)-1)^2}{4\exp(1)}\leq(m:\mathbb R).
\]
Positive matrix size is derived inside the proof.  The stronger count
theorems are retained rather than only their asymptotic consequences.
The direct Zariski quadratic endpoints compose the Zariski count theorems
with \nolinkurl{Sheshadri.ordinary_quadratic_bound} and
\nolinkurl{Sheshadri.symmetric_quadratic_bound}.  They do not
need a Zariski-to-Euclidean closure theorem.

There is also an exact, non-Fermat version.  The theorems
\nolinkurl{ordinary_smooth_degree_bound} and
\nolinkurl{symmetric_smooth_degree_bound} assume $n\geq2$, a nonzero
homogeneous $f$ of positive degree $d$, a literal determinant representation,
and
\[
 x\ne0,\quad f(x)=0\quad\Longrightarrow\quad
 \text{some }\partial_i f(x)\ne0.
\]
They conclude $d(d-1)^{n-2}\leq B(m,n)$ or
$d(d-1)^{n-2}\leq B_{\mathrm{sym}}(m,n)$.  The smoothness and degree
hypotheses are proved for $f_n$ before applying these results.

\subsection{The closure definitions are literal}
\label{app:lean-closure}

Let $S$ be a finite set of exponent vectors.  The coefficient space used in
\nolinkurl{BorderRepresentation.lean} is the finite-dimensional complex
space $\CC^S$, equipped with its usual topology.  A coefficient vector $c$
denotes the polynomial $\sum_{\alpha\in S}c_\alpha X^\alpha$.  The ordinary
locus consists of exactly those vectors whose polynomial is a size-$m$
affine determinant; the symmetric locus additionally requires symmetric
coefficient matrices.  Thus
\texttt{\detokenize{OrdinaryFiniteSupportBorder n m}} means that there exists a finite
$S$ and a vector representing $f_n$ that belongs to the Euclidean closure
of this actual locus.  No family or specialization assertion is part of
this definition.

The canonical \nolinkurl{AtSize} predicates use
\texttt{\detokenize{coefficientSupport n m}}.  The development proves that this
finite set contains the support of $f_n$ and of every size-$m$ affine
determinant.  The theorems \nolinkurl{ordinaryLocus_eq_range} and
\nolinkurl{symmetricLocus_eq_range} identify the loci with the coefficient
images of all such pencils.  Consequently, the canonical support introduces
no unproved restriction on the determinant locus.

The Zariski predicates, defined in \nolinkurl{ZariskiBridge.lean}, use the
common zeros of all polynomials vanishing on the actual coefficient locus.
The proof needed for the Euclidean endpoints is only
\[
 \overline{T}^{\,\mathrm{Euclidean}}
 \ \subseteq\ \overline{T}^{\,\mathrm{Zariski}}.
\]
It is proved using continuity of polynomial evaluation.  The finite-support
coefficient reindexing required to apply this inclusion is also proved.
The final argument then constructs its algebraic family directly from
Zariski membership.  It neither assumes that these two closures agree for
arbitrary sets nor uses their reverse inclusion.  An older conditional
interface named \nolinkurl{UniversalPolynomialImageClosure} is retained in
the source, but it is not an input to any final endpoint listed above.

\subsection{A definition-level dictionary}
\label{app:lean-dictionary}

The formal objects below are concrete algebraic objects, rather than names
for desired numerical properties.  Local module names in this appendix
refer to files in the \nolinkurl{Sheshadri/} directory.

\begingroup\small
\begin{longtable}{@{}>{\raggedright\arraybackslash}p{0.30\textwidth}>{\raggedright\arraybackslash}p{0.64\textwidth}@{}}
\toprule
Mathematical object & Formal realization\\
\midrule
\endhead
Affine determinant and target &
\nolinkurl{PolynomialModel.affinePencil},
\nolinkurl{determinant}, and \nolinkurl{fermat}; actual matrices and
multivariate polynomials.\\[4pt]
Polar section &
\nolinkurl{SmoothGaussSectionLength.polarChartIdeal}; the ideal generated
by $f$, the selected derivative combinations, and $a\cdot x-1$.\\[4pt]
Local intersection algebra &
\nolinkurl{IsolatedLocalMultiplicity.LocalIntersectionRing}; the actual
polynomial quotient localized at the evaluation point.\\[4pt]
Multiplicity &
\nolinkurl{IsolatedLocalMultiplicity.multiplicity}; module length of that
local ring over itself, with proved equality to field dimension in the
rational Artinian setting.\\[4pt]
Special Gauss section &
\nolinkurl{ConormalGaussSection.Section}; actual flags and distinct points,
vanishing, finite quotient, Artinian local rings, smoothness, and a proved
length sum.  \nolinkurl{exists_section} constructs the record.\\[4pt]
Border family &
\nolinkurl{DeterminantDvrFamily.Family}; a DVR, residue and generic maps,
a polynomial over the DVR, one generic pencil, a uniformizer, and literal
specialization and determinant identities.\\[4pt]
Open-locus saturation &
\nolinkurl{HomogeneousSystemSaturation.saturation}; the ideal of $q$ such
that, for every displayed $g_i$, some power $g_i^{N_i}q$ belongs to the
original ideal.\\[4pt]
Finite generic section &
\nolinkurl{ActualGenericGaussSection.exists_actual_section}; constructs
flags with prescribed residues and a finite quotient by the derivative
saturation of the actual generic polar equations.\\[4pt]
Incidence multiplicity &
The actual normalized kernel-incidence equation quotient in a valid
projective chart, localized at its evaluation point; identified with the
polar local ring by proved ring equivalences.\\
\bottomrule
\end{longtable}
\endgroup

For clarity, write the special affine polar ideal as
\[
 I_0=\left(f_n,\
 \sum_i\ell_{1i}\partial_i f_n,\ldots,
 \sum_i\ell_{n-2,i}\partial_i f_n,\ a\cdot x-1\right).
\]
It has $n$ displayed equations in $n$ affine variables.  If $p$ is one of
its selected points, its contribution is
\[
 \mu_p=\operatorname{length}_{(\CC[x]/I_0)_p}
                  (\CC[x]/I_0)_p,
\]
not an assigned weight and not merely the number one.  This distinction
is essential when the polar section is nonreduced.

\subsection{Step 1: construct the special Gauss section and its degree}
\label{app:lean-special}

The construction begins with a homogeneous polynomial that is smooth away
from the cone vertex.  \nolinkurl{SmoothGaussPolars} proves the relevant
Euler identity and gradient basepoint-freeness.  The modules
\nolinkurl{GenericGaussPolars}, \nolinkurl{IrrelevantSaturation}, and
\nolinkurl{GenericGaussPolarSequence} choose successive polar combinations
by avoiding finitely many associated primes.  Irrelevant saturation preserves
sufficiently high graded pieces, while the chosen equations act regularly
on the appropriate saturated quotients.  Generic regularity is thus the
output of an existence proof, not a premise supplied by the caller.

\nolinkurl{SmoothGaussDegree} uses exact sequences of actual graded
quotients to calculate the eventual Hilbert value after the $n-2$ polar
cuts.  It obtains $d(d-1)^{n-2}$.  The construction then chooses a linear
form avoiding the resulting projective support.  Setting this form to one
gives an actual finite-dimensional affine algebra, as proved in
\nolinkurl{SmoothGaussSection} and \nolinkurl{SmoothGaussSectionLength}.

The modules \nolinkurl{ArtinianLocalDecomposition},
\nolinkurl{FiniteQuotientPointFamily}, and
\nolinkurl{SmoothGaussPointFamily} identify the dimension of this finite
algebra with the sum of the full local lengths at its distinct support
points.  The resulting theorem
\nolinkurl{ConormalGaussSection.exists_section} supplies all the data in
the \nolinkurl{Section} record and the identity
\[
 \sum_p\mu_p=d(d-1)^{n-2}.
\]
Specializing to the actual Fermat polynomial gives $\sum_p\mu_p=C(n)$.
The endpoint proof invokes this constructor and discharges its hypotheses;
it does not assume the existence or the numerical degree of such a section.

The chart condition $a\cdot x=1$ is chosen after the homogeneous polar
section.  The proof does not fix an original Fermat coordinate to one and
then explicitly solve the resulting polar equations to obtain the lower
count.  The degree comes from the regular-cut Hilbert calculation.

\subsection{Step 2: construct a genuine DVR family from the border}
\label{app:lean-family}

The border argument starts with the coefficient-space membership of
Section~\ref{app:lean-closure}.  \nolinkurl{AlgebraicCurveSelection} and
\nolinkurl{CurveThroughBoundary} construct a curve through the boundary
point of the polynomial-image closure.  \nolinkurl{NormalGoingDown},
\nolinkurl{CurveNormalization}, and \nolinkurl{FiniteCurveExtension}
construct the normalized branch and the finite extension needed to carry
a coherent generic preimage.  Localizing at the selected point supplies
a discrete valuation ring $D$, a residue map $D\to\CC$, and its actual
fraction field $L$.

\nolinkurl{GenericPolynomialLift} and \nolinkurl{AlgebraicBorderFamily}
produce one tuple of generic matrix parameters.  They do not choose
unrelated representations at separate fibres.  Pole clearing in
\nolinkurl{DvrRationalEntries} uses one common power of a uniformizer.
\nolinkurl{DeterminantDvrFamily} then reconstructs a polynomial
$F\in D[x_1,\ldots,x_n]$ and a pencil $A$ over $L$ satisfying the actual
polynomial identities
\[
 F\bmod\pi=f_n,\qquad F\otimes_D L=\det A(x).
\]
The theorem \nolinkurl{Family.isFractionRing} verifies that the field used
by the pencil is the fraction field of the constructed DVR.  Symmetry is
retained in the symmetric construction.

The final family-selection results are
\nolinkurl{ordinary_border_exact_or_family} and
\nolinkurl{symmetric_border_exact_or_family}.  They return either an exact
representation, to which the exact theorem applies, or the family just
described.  An arbitrary algebraic curve is never replaced by an assumed
parametrization over $\CC[t]$, and the generic field $L$ is never identified
with $\CC$.

\subsection{Step 3: prove that a compatible generic section is finite}
\label{app:lean-generic}

It would be incorrect to infer finite generic fibre merely from finite
special fibre: components can escape the special fibre.  The development
therefore proves generic finiteness separately, using the actual family.

\nolinkurl{DeterminantDvrHomogeneous} homogenizes $F$ in an added variable
and proves the dehomogenization and derivative identities.  Its theorem
\nolinkurl{generic_hypersurface_hilbert} bounds the degree of the initial
Hilbert polynomial by $n-1$.  The saturated generic construction works on
the open locus where the added coordinate and at least one spatial
derivative are nonzero.

The prescribed special flags cannot simply be replaced by unrelated
generic flags.  \nolinkurl{RegularFlagsWithResidue.exists_regular_coefficients}
chooses coefficients in $D$ with exactly the required residues, while
avoiding the finitely many associated primes of the generic quotients.
The choices are perturbations by uniformizer multiples.  The modules
\nolinkurl{PrescribedResidueGaussSequence} and
\nolinkurl{FinitePrescribedGaussSection} apply this construction to all
$n-2$ polar cuts and the final primal cut.  Normalization of the primal
equation divides only by a proved unit of $D$ and preserves its prescribed
residue.

\nolinkurl{HomogeneousHilbertDegree} proves the degree decrease under a
regular homogeneous cut and the required monotonicity under enlargement
of the ideal.  \nolinkurl{SaturatedSectionHilbert} applies these results
through the actual cut sequence.  After $n-1$ cuts the Hilbert function is
eventually constant.  Regularity of the added coordinate then proves that
the corresponding affine chart algebra is finite-dimensional.

Dehomogenization identifies an actual ideal $J\subseteq L[x]$.  If
\[
 I_\eta=\left(\det A,\
 \sum_i\widetilde\ell_{1i}\partial_i\det A,\ldots,
 \sum_i\widetilde\ell_{n-2,i}\partial_i\det A,\
 \widetilde a\cdot x-1\right),
\]
then \nolinkurl{ActualGenericGaussSection.exists_actual_section} proves
\begin{gather*}
 \dim_L L[x]/J<\infty,\\
 I_\eta\subseteq J\subseteq
 \operatorname{Sat}_{\partial\det A}(I_\eta),\qquad
 J=\operatorname{Sat}_{\partial\det A}(I_\eta).
\end{gather*}
Here $q\in\operatorname{Sat}_{g}(I)$ means
$\forall i\,\exists N_i\geq0:\ g_i^{N_i}q\in I$.  This operation removes
components supported on the excluded closed locus; it does not take a
radical.  Where one $g_i$ is a unit, the localized saturated and original
ideals are proved equal.

Finally, \nolinkurl{SaturatedChartSmoothness} constructs a derivative
combination whose class in the full finite quotient $L[x]/J$ is a unit.
Regularity gives injectivity of multiplication; finite dimension gives
surjectivity and hence the inverse.  Therefore all geometric support
points of $J$ are smooth points of the determinant hypersurface, including
after passage to an algebraic closure.  The full quotient by $I_\eta$
need not be finite: excess components outside this smooth section are
allowed throughout.

\subsection{Step 4: preserve singular local multiplicities}
\label{app:lean-persistence}

The lower comparison must preserve nilpotents, not merely show that the
selected special points are in a limiting support.  This is the role of
\nolinkurl{DVRRootPersistence}.

At a selected special point, the actual local polar algebra is Artinian
and is given by a square equation system.  The regular-sequence machinery
in \nolinkurl{CoordinatePowerRegular},
\nolinkurl{LocalCompleteIntersection}, and
\nolinkurl{IsolatedCompleteIntersection} derives regularity of these
equations in the local polynomial ring.  It does not require a nonzero
Jacobian determinant.  The total local equation quotient then has regular
uniformizer action.  Cancelling generic denominators shows that contraction
from the generic fibre annihilates the entire original special local
algebra.

\nolinkurl{DVRSaturatedPersistence} handles the additional saturation in
$J$.  At a selected special point, one spatial derivative has nonzero
value, and consequently is a unit in its full Artinian local ring.  Power
membership in the saturation, followed by cancellation of this unit, shows
that every contracted element of $J$ still annihilates that ring.
\nolinkurl{DVRSaturatedGaussPersistence} applies this result to the actual
family and the residue-preserving flags.  Its theorem
\nolinkurl{section_saturated_contraction_local_eq} proves equality of the
corresponding localized ideals.  Thus this step retains local scheme
structure and local length.

Let $T=J\cap D[x]$ and let $T_0$ be its residue ideal.  Simultaneous
interpolation onto the full special local algebras gives
\[
 \sum_p\mu_p\ \leq\ \dim_{\CC}\CC[x]/T_0.
\]
The quotient $D[x]/T$ embeds in its generic fibre and is torsion-free over
$D$.  \nolinkurl{DVRFibreDimension} proves the residue-dimension inequality
\[
 \dim_{\CC}\CC[x]/T_0\ \leq\ \dim_L L[x]/J.
\]
It does not assume finite generation of the whole horizontal algebra over
$D$; the already proved finite generic dimension is used where needed.
\nolinkurl{DVRSectionLength.selected_lengths_le_generic} combines these
facts, and \nolinkurl{DVRGaussLengthBound.family_section_degree_le}
instantiates them to obtain
\begin{equation}
 \label{eq:lean-special-generic}
 n(n-1)^{n-2}\ \leq\ \dim_L L[x]/J.
\end{equation}

\subsection{Step 5: B\'ezout with isolated singular zeros and excess}
\label{app:lean-bezout}

The precise B\'ezout statement matters.  The development proves an
isolated-intersection inequality; it does not assert that an arbitrary
nonproper intersection has the expected total length.  The theorem
\nolinkurl{FieldSingularBezoutWithExcess.isolated_multiplicity_bound}
is over a characteristic-zero field and has the following inputs:
\begin{enumerate}[leftmargin=*]
\item Three finite nonempty homogeneous coordinate blocks of projective
dimensions $a,b,c$, and exactly $a+b+c$ actual equations.
\item The displayed multihomogeneity of each equation, with multidegree
$(d_{e,1},d_{e,2},d_{e,3})$.
\item A finite family of common zeros, each with a specified nonzero chart
coordinate in every block, distinct modulo the independent projective
block scalings.
\item Artinianness of each selected point's actual localized chart
equation ring.
\end{enumerate}
Its conclusion is
\begin{equation}
 \label{eq:lean-bezout}
 \sum_p\operatorname{length}(\mathcal A_p)
 \ \leq\
 [H_1^aH_2^bH_3^c]
 \prod_e(d_{e,1}H_1+d_{e,2}H_2+d_{e,3}H_3).
\end{equation}
The family of selected zeros need not exhaust the zero set.  Other
components may have positive dimension.  The selected roots need not be
simple or reduced.  A companion theorem,
\nolinkurl{isolated_multiplicity_bound_of_isOpen}, derives Artinianness
from literal isolatedness expressed as an open singleton in the actual
chart spectrum.

The proof of~\eqref{eq:lean-bezout} is itself part of the formalization.
It interpolates the given equations with a split comparison system of the
same multidegrees, contracts the generic ideal without radicalizing it,
and proves flatness and conservation of actual finite multigraded pieces.
Local complete-intersection persistence retains all selected local
lengths at the original endpoint.  At the comparison endpoint, the split
equation quotient surjects onto the contracted-fibre quotient, giving the
upper-bound inequality in the correct direction.

The split-system estimate uses a Hilbert-function inequality for products,
not an unproved regular-sequence equality in the entire homogeneous
coordinate ring.  Expanding products gives finitely many branches.
Branches selecting too many linear equations in a block vanish in
sufficiently positive degree; branches meeting exactly the dimension
quotas contribute at most one.  The number of quota choices is the
coefficient on the right of~\eqref{eq:lean-bezout}.  All required degree
thresholds and componentwise shifts are proved.  This approach allows
irrelevant and excess components without identifying point counts with
the degree of an arbitrary homogeneous quotient.

The main dependency modules for this part are
\nolinkurl{FieldSingularBezoutDeformationBound},
\nolinkurl{FieldBezoutSplitHilbert},
\nolinkurl{FieldBezoutVariableTransport}, and
\nolinkurl{FieldSingularBezoutWithExcess}, together with their local
persistence and multigraded algebra dependencies.  The comparison split
system is a general B\'ezout upper-bound construction; it is unrelated to
a fixed-coordinate explicit count for the Fermat polynomial.

\subsection{Step 6: apply B\'ezout to the actual determinant incidence}
\label{app:lean-incidence}

At a smooth determinant zero, the proof constructs normalized kernel
vectors and valid cofactor charts.  In the ordinary model it also
constructs one reduction matrix valid simultaneously for the finite
selected family.  The Schur-complement calculations are equalities of
ideals and actual quotient/local-ring equivalences, not only equalities
of their zero sets.  The symmetric model uses its native one-kernel-vector
incidence rather than discarding the symmetry condition.

\nolinkurl{FieldPolarMultiplicityBound} applies the singular isolated-zero
B\'ezout theorem to these literal incidence equations.  It proves the
equation count, multidegrees, projective distinctness and Artinian local
conditions from the constructed polar data.  The local equivalences
transfer the full intersection lengths.  The resulting incidence
coefficients are precisely $B(m,n)$ and $B_{\mathrm{sym}}(m,n)$.

The generic field of the border family need not be algebraically closed.
\nolinkurl{ArbitraryFieldFinitePolarBound} therefore extends the actual
finite algebra $L[x]/J$ to an algebraic closure $\overline L$.  Explicit
tensor/quotient equivalences prove
\[
 \dim_{\overline L}\bigl(\overline L\otimes_L L[x]/J\bigr)
 =\dim_L L[x]/J.
\]
The unit derivative combination remains a unit after extension, so every
support point over $\overline L$ is smooth.  The proved saturation
inclusions suffice to identify each localized $J$-algebra with the
original polar algebra.  The argument uses the required inclusion under
mapping an ideal; it does not assume that saturation commutes with every
base change.

\nolinkurl{FieldFinitePolarBound} constructs the complete point family
of the finite quotient and sums its local lengths.  Consequently
\nolinkurl{ArbitraryFieldFinitePolarBound.ordinary_finite_saturated_polar_bound}
and its symmetric counterpart give
\begin{equation}
 \label{eq:lean-generic-upper}
 \dim_L L[x]/J\leq B(m,n),
 \qquad\text{or}\qquad
 \dim_L L[x]/J\leq B_{\mathrm{sym}}(m,n).
\end{equation}
No generic homogeneity of the original affine determinant polynomial and
no finite-dimensionality of its entire unsaturated polar quotient is
assumed here.

\subsection{Step 7: assemble the endpoints and the constants}
\label{app:lean-assembly}

The family-level theorem
\nolinkurl{ConormalBorderMultiplicity.ordinary_family_count} constructs
the special section and the compatible finite generic section, applies
\eqref{eq:lean-special-generic} and~\eqref{eq:lean-generic-upper}, and
concludes $C(n)\leq B(m,n)$.  The symmetric theorem performs the same
assembly with the symmetric incidence.  The Zariski endpoints invoke
the exact-or-family constructor, and the Euclidean endpoints invoke the
proved one-way closure transfer.  Thus the intermediate family and section
records are outputs of proved constructors, not undeclared obligations
of the endpoint user.

\nolinkurl{Bezout.lean} proves the coefficient identities and the
Vandermonde estimate.  \nolinkurl{Numerical.lean} proves
\[
 B(m,n)\leq\left(\frac{4em}{n-1}\right)^{n-1},
 \qquad
 B_{\mathrm{sym}}(m,n)
 \leq\left(\frac{2em}{n-1}\right)^{n-1}.
\]
Combining these with $(n-1)^{n-1}\leq C(n)$ and the positivity following
from $n\geq3$ yields the two quadratic bounds.  All casts between natural
numbers and real numbers, the exponential constants, and the power
comparison are checked by Lean.

The exact proof is shorter: it constructs the special section, identifies
its actual local rings with the kernel incidence, and applies B\'ezout
directly.  It does not need the border family or specialization comparison.

\subsection{What the formalization establishes, and what it does not}
\label{app:lean-scope}

The verified numerical chain for the border theorem is
\[
 \underbrace{n(n-1)^{n-2}}_{\text{constructed special local lengths}}
 \ \leq\
 \underbrace{\dim_L L[x]/J}_{\text{constructed generic section}}
 \ \leq\
 \underbrace{B(m,n)}_{\text{proved incidence B\'ezout coefficient}},
\]
with the symmetric replacement in the last term.  This is a local
Gauss-section implementation of the conormal/Gauss specialization
argument.  It establishes the stated exact and border lower-bound
endpoints without assuming a global projective cycle/intersection
certificate.

The written proof also distinguishes numerical
constancy of fibre multidegrees from rational equivalence of fibres over
an arbitrary base curve: the relevant numerical constancy is obtained
from the bigraded Euler--Hilbert calculation.  Likewise, the cone
comparison is made away from the vertex, and the primal hyperplane is
chosen to avoid the vertex.  No global projective-line-bundle description
through the vertex is needed.  These distinctions are consistent with
the local-section implementation; the end-to-end Lean guarantee remains
the precise endpoint statements in Section~\ref{app:lean-statements}.

The development also contains separately proved results on actual Gauss
families, flatness, conormal containment, grading, Hilbert functions, and
projective charts.  It does not claim a separately complete global Chow
ring library, the full global projective fibre-cycle assembly of the main
text, or an identification of every concrete local chart ring with a
stalk on a product of Mathlib \nolinkurl{Proj} schemes.  Those broader
claims are not needed as hypotheses of the assembled endpoints.

Older modules retain conditional interfaces called
\nolinkurl{GaussDegreeTheory} and \nolinkurl{FamilyCertificate}.  A theorem
accepting one of these records remains conditional.  Importing its module
does not make the record an axiom, but a clean axiom list also does not
prove the record exists.  The final \nolinkurl{Multiplicity} endpoint
modules use the constructed local-section route described above and do
not require those interfaces.

The only affine chart used in the generic construction sets the
\emph{added homogenizing coordinate} $x_0$ to one.  No original coordinate
of $F_n$ is set to one, and no explicit count of polar points is used: the
lower count comes from the Hilbert-function calculation of
Section~\ref{app:lean-special}.

\subsection{What the audits check}
\label{app:lean-audits}

The release's entry-point dependency graph contains 225 local modules.
The principal axiom audit requests 349 reports.  The accompanying
\nolinkurl{BUILD_LOG.txt}, \nolinkurl{AXIOM_AUDIT.txt}, and
\nolinkurl{IMPORT_AUDIT.txt} record successful build and audit results.
The scripts recompute the checks from the delivered source; these counts
are descriptive release metadata rather than mathematical assumptions.
The implementation and source-level review were AI-assisted.  The
semantic audit documents record separate automated source reviews, not
independent human peer review.  The compiler and axiom transcripts are
the reproducible machine-checking evidence.

There are three distinct checks.
\begin{enumerate}[leftmargin=*]
\item \textbf{Compilation.}  Lean elaborates the local development against
the pinned dependencies and checks the resulting proof terms.  A build
failure is not converted into a successful verification claim.
\item \textbf{Foundational dependencies.}
\nolinkurl{AxiomAudit.lean} runs \texttt{\detokenize{#print axioms}} on the principal
declarations, including the final endpoints.  The only reported axioms
are Lean's standard \nolinkurl{propext}, \nolinkurl{Classical.choice}, and
\nolinkurl{Quot.sound}.  \nolinkurl{audit_axioms.py} checks that every
requested report is present and rejects unexpected dependencies.
\nolinkurl{audit_imports.py} checks the local entry-point import graph for
missing files, custom axioms, unfinished proofs, unsafe shortcuts and the
excluded alternative explicit-slice modules.
\item \textbf{Mathematical statement and applicability.}
\nolinkurl{END_TO_END_HYPOTHESES_AUDIT.txt} reviews the literal definitions,
the final theorem types, and how the proof supplies every constructor's
applicability conditions.  This check is necessary because a theorem of
the form ``if the desired count bound holds, then the desired conclusion
holds'' can compile with only standard foundational axioms while proving
little about the intended geometric problem.
\end{enumerate}

These checks address different questions.  A source scan is not a
substitute for compilation, an axiom report is not a semantic audit, and
the existence of a successful transcript is not a substitute for rebuilding
the sources when independent reproduction is desired.  The trusted base
includes the Lean kernel and the pinned mathematical definitions and
dependencies.  The verification is a proof for all admitted $n,m$, not
numerical testing of finitely many dimensions.

Construction-specific reading guides are
\nolinkurl{GAUSS_DEGREE_CONSTRUCTION_AUDIT.txt},
\nolinkurl{BORDER_DVR_CONSTRUCTION_AUDIT.txt},
\nolinkurl{GENERIC_GAUSS_SECTION_AUDIT.txt},
\nolinkurl{DVR_LOCAL_PERSISTENCE_AUDIT.txt}, and
\nolinkurl{EXACT_MULTIPLICITY_AUDIT.txt}.  Historical audit files are
retained to document earlier stages; their earlier lists of remaining
obligations are superseded by \nolinkurl{ORIGINAL_CONORMAL_SCOPE.txt} and
\nolinkurl{END_TO_END_HYPOTHESES_AUDIT.txt}.

\subsection{Reproducing the build from the release sources}
\label{app:lean-reproduce}

The release is source-complete for its local entry-point graph.  It does
not vendor the Lean compiler, Mathlib, or precompiled Mathlib caches.
A fresh machine therefore needs an Internet connection for the initial
toolchain and dependency download, Git, Python~3, a shell capable of
running the supplied Bash script, and Lean's toolchain manager Elan.
The mathematical toolchain is pinned to
\begin{itemize}[leftmargin=*]
\item Lean \nolinkurl{leanprover/lean4:v4.19.0};
\item Mathlib revision
\nolinkurl{c44e0c8ee63ca166450922a373c7409c5d26b00b}
(Mathlib v4.19.0);
\item the remaining dependency revisions in the delivered
\nolinkurl{lake-manifest.json}.
\end{itemize}
Download and extract the archived release, or clone the repository (both linked at the start of this appendix), and keep its manifest.  From the
directory containing \nolinkurl{lean-toolchain}, \nolinkurl{lakefile.toml},
and \nolinkurl{verify.sh}, run:

\begingroup\small
\begin{verbatim}
elan toolchain install leanprover/lean4:v4.19.0
lake env lean --version
lake exe cache get
./verify.sh
\end{verbatim}
\endgroup

If the extraction did not preserve executable permissions, use
\texttt{\detokenize{bash verify.sh}}.  The script executes, in order:

\begingroup\small
\begin{verbatim}
lake build
lake env lean AxiomAudit.lean | tee AXIOM_AUDIT.txt
python3 audit_axioms.py
python3 audit_imports.py
\end{verbatim}
\endgroup

It stops on failure.  A successful return code is zero.  The expected
audit summaries report all requested declarations, only the standard
foundational axioms, and a complete entry-point source graph with no
prohibited shortcuts.  The Mathlib cache is a build accelerator; the
release supplies source proofs for its own modules.  A reader wishing
to rebuild dependencies from source may omit the cache download and
allow the dependency build to complete instead.

The manifest-preserving build above is the reproduction procedure for
the archived release.  \texttt{\detokenize{lake update}} is useful when deliberately
regenerating dependency metadata, but is not necessary for this procedure
and may change transitive dependency resolutions.  If it is used, retain
a copy of the archived manifest, compare every resolved revision with
that copy, and restore the archived manifest before claiming an exact
release reproduction.  The Mathlib checkout can be checked directly:

\begingroup\small
\begin{verbatim}
git -C .lake/packages/mathlib rev-parse HEAD
\end{verbatim}
\endgroup

It must print the Mathlib commit above.  The release's
\nolinkurl{SOURCE_SHA256SUMS.txt} records the delivered file contents
except the checksum manifest itself.  Verify these checksums on the
untouched extraction before rerunning scripts that regenerate logs.
On a system with the GNU checksum utility, run
\texttt{\detokenize{sha256sum -c SOURCE_SHA256SUMS.txt}} from the directory to which
the manifest paths are relative.

\nolinkurl{BUILD_ENVIRONMENT.txt} describes the original build environment
and the small executable-location compatibility wrapper needed by that
container.  The wrapper changes only the process executable-location
lookup and is not part of the mathematical development.  A normal Lean
installation does not need it.  The compiler and shared-library hashes
are recorded there for inspection.

\subsection{A short independent endpoint inspection}
\label{app:lean-smokecheck}

A reader need not begin by reading every implementation file.  After
building the release, create a new file named
\nolinkurl{EndpointInspection.lean} in the project root with the following
contents.  This file only inspects the already defined statements; it
does not introduce assumptions or replace a proof.

\begingroup\small
\begin{verbatim}
import OriginalConormal

set_option pp.explicit true
set_option pp.universes true

open Sheshadri.ConormalExactMultiplicity in
#check @ordinary_exact_quadratic
open Sheshadri.ConormalExactMultiplicity in
#check @symmetric_exact_quadratic

open Sheshadri.ConormalBorderMultiplicity in
#check @ordinary_border_quadratic
open Sheshadri.ConormalBorderMultiplicity in
#check @symmetric_border_quadratic
open Sheshadri.ConormalBorderMultiplicity in
#check @ordinary_borderAtSize_quadratic
open Sheshadri.ConormalBorderMultiplicity in
#check @symmetric_borderAtSize_quadratic
open Sheshadri.ConormalBorderMultiplicity in
#check @ordinary_zariski_count
open Sheshadri.ConormalBorderMultiplicity in
#check @symmetric_zariski_count
open Sheshadri.ConormalBorderMultiplicity in
#check @ordinary_zariski_quadratic
open Sheshadri.ConormalBorderMultiplicity in
#check @symmetric_zariski_quadratic

open Sheshadri.ConormalBorderMultiplicity in
#print axioms ordinary_border_quadratic
open Sheshadri.ConormalBorderMultiplicity in
#print axioms symmetric_border_quadratic

#print Sheshadri.PolynomialModel.ExactRepresentation
#print Sheshadri.PolynomialModel.SymmetricExactRepresentation
open Sheshadri.BorderRepresentation in
#print OrdinaryFiniteSupportBorder
open Sheshadri.BorderRepresentation in
#print OrdinaryBorder
open Sheshadri.BorderRepresentation in
#print ordinaryLocus
open Sheshadri.ZariskiBridge in
#print OrdinaryZariskiBorder
open Sheshadri.ZariskiBridge in
#print zariskiClosure
\end{verbatim}
\endgroup

Run \texttt{\detokenize{lake env lean EndpointInspection.lean}}.  The output makes
all implicit parameters visible.  In particular, the final exact and
border statements must have only the natural numbers $n,m$, the condition
$3\leq n$, and the literal representation or border predicate.  The
ordinary and symmetric endpoint axiom lists must contain no custom axiom
or \nolinkurl{sorryAx}.  Repeat the definition inspection with the symmetric
border names if desired; their extra condition is literal matrix symmetry.

For a mathematical review, a useful order is: inspect the definitions and
final types; read the two endpoint modules; follow
\nolinkurl{exists_section} and \nolinkurl{exists_actual_section}; inspect
\nolinkurl{family_section_degree_le}; then inspect the arbitrary-field
finite polar bound and the singular B\'ezout statement.  At each step,
check both the intermediate theorem's assumptions and the constructor
that supplies them.  This follows the actual proof dependency chain and
separates a reusable conditional lemma from an assembled unconditional
endpoint.

\endgroup
